%% ma: L12-B15-0002
%% lop: 12
%% chuong: 5
%% bai: 15
%% dang: 12.15.2
%% loai: TN4
%% muc: NB
%% dap_an: "C"
%% nguon: "(NBV)-ĐỀ SỐ 1-MỖI NGÀY 1 ĐỀ THI 2025.pdf (D01, MỖI NGÀY 1 ĐỀ - Đề số 1), Phần I Câu 7"
%% kien_thuc: "Bài 15 (phương trình chính tắc của đường thẳng, vectơ chỉ phương)"
%% trang_thai: chinh-thuc
\begin{ex}
Trong không gian $Oxyz$, cho đường thẳng $d:\dfrac{x-1}{4}=\dfrac{-y}{2}=\dfrac{z+2}{-6}$. Vectơ nào dưới đây là một vectơ chỉ phương của đường thẳng $d$?
\choice{$\vec u_2(2;-1;3)$}{$\vec u_1(4;2;-6)$}{\True $\vec u_3(-2;1;3)$}{$\vec u_4(1;0;2)$}
\loigiai{Viết lại $d:\dfrac{x-1}{4}=\dfrac{y}{-2}=\dfrac{z+2}{-6}$, nên $d$ có vectơ chỉ phương $(4;-2;-6)=-2\cdot(-2;1;3)$. Vậy $\vec u_3(-2;1;3)$ là một vectơ chỉ phương.}
\end{ex}
